\chapter{Acteurs et Rôles}

\section{Gouvernance du Projet}

Conformément aux principes de gestion de projet (Chapitre 2 du cours), nous avons identifié les parties prenantes selon le modèle classique \textbf{MOA / MOE / Chef de Projet}.

\subsection{Maîtrise d'Ouvrage (MOA) --- Le Client}

La MOA représente l'\textbf{équipe Marketing} de notre plateforme e-commerce fictive. Son rôle est de définir le besoin :

\begin{table}[H]
\centering
\begin{tabularx}{\textwidth}{L{4cm} X}
\toprule
\textbf{\color{primarycolor}Aspect} & \textbf{\color{primarycolor}Détail} \\
\midrule
\textbf{Qui ?} & Directeur Marketing de la plateforme e-commerce \\
\addlinespace[0.1cm]
\textbf{Besoin principal} & Comprendre le comportement des clients en temps réel \\
\addlinespace[0.1cm]
\textbf{Objectifs} & Réduire le taux d'abandon de panier, personnaliser l'expérience \\
\addlinespace[0.1cm]
\textbf{Exigences} & Voir les clics en temps réel, être alerté des abandons de panier, obtenir des statistiques de vente par produit \\
\addlinespace[0.1cm]
\textbf{Validation} & Vérifier que les données affichées correspondent à la réalité \\
\bottomrule
\end{tabularx}
\caption{Fiche de la Maîtrise d'Ouvrage}
\end{table}

\subsection{Maîtrise d'Œuvre (MOE) --- L'Équipe Technique}

La MOE est notre \textbf{équipe de développement}. Elle est chargée de concevoir et d'implémenter la solution technique répondant aux exigences de la MOA.

\begin{table}[H]
\centering
\begin{tabularx}{\textwidth}{L{4cm} X}
\toprule
\textbf{\color{primarycolor}Aspect} & \textbf{\color{primarycolor}Détail} \\
\midrule
\textbf{Qui ?} & L'équipe d'étudiants développeurs \\
\addlinespace[0.1cm]
\textbf{Rôle} & Concevoir et réaliser l'architecture technique \\
\addlinespace[0.1cm]
\textbf{Responsabilités} & Déployer l'infrastructure (Docker, Kafka, Flink), développer les topologies de traitement, configurer la rétention \\
\addlinespace[0.1cm]
\textbf{Contraintes} & Technologies imposées : Kafka, Flink, Docker \\
\bottomrule
\end{tabularx}
\caption{Fiche de la Maîtrise d'Œuvre}
\end{table}

\subsection{Chef de Projet}

Le chef de projet fait le lien entre la MOA et la MOE. Il s'assure que la MOE ne construit pas une solution trop complexe (<<~usine à gaz~>>) et que la MOA comprend les contraintes techniques, notamment les enjeux de stockage et de rétention des données.

\begin{table}[H]
\centering
\begin{tabularx}{\textwidth}{L{4cm} X}
\toprule
\textbf{\color{primarycolor}Aspect} & \textbf{\color{primarycolor}Détail} \\
\midrule
\textbf{Rôle} & Coordonner la MOA (Marketing) et la MOE (Dev) \\
\addlinespace[0.1cm]
\textbf{Missions} & Planifier les sprints, gérer les risques, assurer la communication \\
\addlinespace[0.1cm]
\textbf{Outils} & Taiga (Kanban), GanttProject, Docker Desktop \\
\bottomrule
\end{tabularx}
\caption{Fiche du Chef de Projet}
\end{table}

\section{Matrice RACI}

La matrice RACI définit clairement les responsabilités pour chaque activité du projet :

\begin{table}[H]
\centering
\begin{tabularx}{\textwidth}{X C{2.5cm} C{2.5cm} C{2.5cm}}
\toprule
\textbf{\color{primarycolor}Activité} &
\textbf{\color{primarycolor}MOA} &
\textbf{\color{primarycolor}Chef de Projet} &
\textbf{\color{primarycolor}MOE} \\
\midrule
Définir les besoins & \textbf{R} & C & I \\
Choisir l'architecture & C & \textbf{A} & \textbf{R} \\
Planifier les sprints & I & \textbf{R} & C \\
Déployer l'infrastructure & I & I & \textbf{R} \\
Développer les topologies & I & C & \textbf{R} \\
Configurer la rétention & I & I & \textbf{R} \\
Tester la cohérence (Replay) & C & \textbf{A} & \textbf{R} \\
Valider les résultats (Recette) & \textbf{R} & A & I \\
Rédiger le rapport & C & \textbf{R} & C \\
Présenter le projet & \textbf{R} & \textbf{R} & \textbf{R} \\
\bottomrule
\end{tabularx}
\caption{Matrice RACI --- \textit{R = Responsable, A = Approbateur, C = Consulté, I = Informé}}
\end{table}

\section{Communication entre Acteurs}

\subsection{Réunions}

La communication entre les acteurs est facilitée par des réunions régulières :
\begin{itemize}
    \item \textbf{Daily Standup} (informel) : point quotidien entre les membres de l'équipe~;
    \item \textbf{Sprint Review} : à chaque fin de sprint, démonstration des livrables à la MOA~;
    \item \textbf{Sprint Retrospective} : retour d'expérience pour améliorer le processus.
\end{itemize}

\subsection{Outils de Communication}

\begin{table}[H]
\centering
\begin{tabularx}{\textwidth}{L{4cm} X}
\toprule
\textbf{\color{primarycolor}Outil} & \textbf{\color{primarycolor}Usage} \\
\midrule
\textbf{Taiga} & Kanban board --- suivi des tâches et User Stories \\
\textbf{GanttProject} & Planification temporelle et chemin critique \\
\textbf{WhatsApp / Discord} & Communication instantanée au quotidien \\
\textbf{Git} & Gestion du code source et versioning \\
\bottomrule
\end{tabularx}
\caption{Outils de communication du projet}
\end{table}
